\DTMsavedate{mydate}{2022-06-30}
\maketitle{Electromagnetic Theory}{Electronic Engineering}{\DTMusedate{mydate}}

% ----------------------------------------------------- %
% COURSE INFO
\subsection*{Course Information}\label{course:3201}
\vspace{0ex}%
\noindent\parbox{0.5\textwidth}{%
    \noindent\begin{tabular}{@{}ll}
                 \textsf{Course:}          & 	Electromagnetic theory \\
                 \textsf{Course Code:}         & 	TEE 3201    \\
                 \textsf{Credits:}    & 	10
    \end{tabular}}
\vspace{2ex}\hrule\vspace{2ex}

% ----------------------------------------------------- %
% INSTRUCTOR INFO
\subsection*{Lecturer Information}
\noindent\parbox{0.5\textwidth}{%
    \noindent\begin{tabular}{@{}ll}
                 \textsf{Name:}         &    Innocent Muchingami           \\
% \textsf{Phone:}        &    \phone          \\
\textsf{Email:}        &    \href{mailto:innocent.muchingami@nust.ac.zw}{innocent.muchingami@nust.ac.zw}    \\
\textsf{Qualifications:}        & PhD in Groundwater Contaminant Transport (UWC), MSc in Geophysics \\
 &                                (NUST), PGD in Higher Education (NUST), BSc (Hons) in Applied Physics \\
&                                 (NUST)

    \end{tabular}}
\vspace{2ex}\hrule\vspace{2ex}


% ----------------------------------------------------- %
% COURSE SYNOPSIS
\subsection*{Course Synopsis}
The module covers Maxwell's equations i.e.\ Laplace and Poisson equations and their solution, Boundary conditions, Plane waves in a perfect dielectric, propagation in imperfect dielectric, Propagation in  imperfect conductors, skin effect, Generalized wave equation, field distributions in rectangular waveguide, Radiation field, dipoles, radiation resistance, impedance, mutual impedance and linear arrays.
\vspace{2ex}\hrule\vspace{2ex}

% ----------------------------------------------------- %
% LEARNING OBJECTIVES
\subsection*{Learning Objectives}
\begin{outline}
    \1 Understand the importance of Vector Analysis in electromagnetism.

    \1 Describe and understand the basic concepts of electrostatics and magneto-statics.

    \1 Master the basic concepts of electrodynamics.

    \1 Understand the laws on conservation as applied to electrodynamics.

    \1 Apply the vector analysis in understanding electromagnetic wave propagation.

    \1 Mathematically solve problems related to electromagnetism using methods of Vector Analysis.

\end{outline}
\vspace{2ex}\hrule\vspace{2ex}

% ----------------------------------------------------- %
% COURSE TOPICS
\subsection*{Topics}
\begin{outline}
    \1 Vector Analysis
    \1 Electrostatics
    \1 Electric field in Matter
    \1 Magneto-statics
    \1 Magnetic fields in Matter
    \1 Electrodynamics
    \1 Conservation laws
    \1 Electromagnetic waves
\end{outline}
\vspace{2ex} \hrule\vspace{2ex}

% ----------------------------------------------------- %
% Students Assenssment
\subsection*{Assessment of Students}
Two assignments and two tests as course work.
Three hour final examination at the end of the semester.
\vspace{2ex}\hrule\vspace{2ex}

% ----------------------------------------------------- %
% Grade Evaluation
\subsection*{Course Evaluation and Grading Policies}
Grades are based on:
Continuous work (\qty{25}{\percent}),
Final Examination (\qty{75}{\percent})
\vspace{2ex}\hrule\vspace{2ex}
% ----------------------------------------------------- %
% EQUIPMENT
% https://sites.udel.edu/computing-purchases/personal-specs/
% \subsection*{Essential Equipment}

% \vspace{2ex}\hrule\vspace{2ex}

% Policies and Other information
\subsection*{Policies and Other Information}
% Key Text or some Books
\subsection*{Key Text}

\begin{itemize}
    \item Electromagnetics by J. D. Kraus and K. R. Carver
    \item Classical electrodynamics by Jackson
    \item Classical electromagnetic theory by Jack Vanderlinde
    \item Introduction to electrodynamics by D. J. Griffiths
\end{itemize}
\vspace{2ex}\hrule
\vspace{2ex}\hrule\vspace{2ex}